\documentclass[10pt,a4paper]{amsart}
\usepackage[margin=19mm]{geometry}
\usepackage{amsmath,amssymb,mathtools}
\usepackage[scr=rsfs,cal=euler]{mathalfa}
\usepackage{xfrac}
\usepackage[unicode,hidelinks,bookmarks=false]{hyperref}
\newcommand{\ie}{i.e.\ }
\newcommand{\cf}{cf.\ }
\newcommand{\resp}{respectively\ }
\newcommand{\Subsection}{\subsection}
\providecommand{\nobreakdash}{\nobreak\hbox{-}}
\setlength{\emergencystretch}{2em}
\multlinegap0pt
\usepackage{bbm}
\usepackage{ulem}
\usepackage{color,framed}
\definecolor{darkgreen}{rgb}{0,0.8,0.2}
\hypersetup{
    colorlinks=true,
    citecolor=darkgreen,
    linkcolor=blue,
    urlcolor=magenta
}
\usepackage{lipsum}
\usepackage{slashed}
\usepackage{url}
\usepackage{bbm}
\usepackage{chngcntr}
\usepackage{mathrsfs}
\usepackage{braket}
\usepackage{tikz}
\usetikzlibrary{decorations.pathreplacing,calc}
\definecolor{myred}{HTML}{a82a2a}
\def\eea{\end{eqnarray}}
\def\rightar{\rightarrow}
\def\lefta{\lambda}
\def\al{\alpha}
\def\prl{\parallel}
\def\ga{\gamma}
\def\om{\omega}
\def\psit{\tilde{\psi}}
\def\Tr{ {\rm Tr} }
\def\<{\langle}
\def\>{\rangle}
\def\mx{\mathrm{x}}
\def\mm{\mathrm{m}}
\def\bG{\mathbb{G}}
\def\bZ{\mathbb{Z}}
\def\bR{\mathbb{R}}
\def\bC{\mathbb{C}}
\def\bE{\mathbb{E}}
\def\cL{\mathcal{L}}
\def\cG{\mathcal{G}}
\def\cA{\mathcal{A}}
\def\cO{\mathcal{O}}
\def\cC{\mathcal{C}}
\def\cH{\mathcal{H}}
\def\cT{\mathcal{T}}
\def\cL{\mathcal{L}}
\def\uX{\underline{X}}
\def\uY{\underline{Y}}
\def\uZ{\underline{Z}}
\newtheorem{thm}{Theorem}
\newtheorem{qu}{Question}
\newtheorem{lemma}{Lemma}
\newtheorem{prop}{Proposition}
\newtheorem{conj}{Conjecture}
\newtheorem{cor}{Corrollary}
\newtheorem{defin}{Definition}
\newtheorem{q}{Question}
\usepackage{color}
\usepackage{tikz-cd}
\newcommand{\de}[1]{\textcolor{red}{DE: #1}}
\newcommand{\eqnref}[1]{Eq.~(\rightef{#1})}
\newcommand{\U}{\mathrm{U}}
\usepackage{tikz}
\usepackage{adjustbox}
\def\[#1\]{%
\begin{equation}\begin{gathered}#1\end{gathered}\end{equation}%
}
\def\Sum_#1{
\sum_{\substack{#1}}
}
\begin{document}
\title[Perturbatively Stable Self-Correcting Classical Memory from ]{Perturbatively Stable Self-Correcting Classical Memory from Gauge Averaging}
\author{Ryan Thorngren}
\date{}
\maketitle
\noindent\emph{Source and licence.} Perturbatively Stable Self-Correcting Classical Memory from Gauge Averaging. Version arXiv:2607.20605v1 (CC BY 4.0). This material is distributed under the Creative Commons Attribution 4.0 International licence, http://creativecommons.org/licenses/by/4.0/, as recorded in the arXiv OAI licence metadata for this exact version and shown on the abstract page https://arxiv.org/abs/2607.20605v1. Full rights notice: licenses/arxiv-2607.20605-CC-BY-4.0.txt.\par\smallskip
\noindent\textit{Editorial scope.} Counted unit: the original \texttt{thm} environment \texttt{thmWegnermemoryapp} at source lines 491--500 together with its complete proof at lines 501--565; the delivered document keeps source lines 260--566 so that every referenced label and every citation resolves inside it. Native editable source is the paper's own concatenated TeX; the AMS wrapper is editorial formatting only. No publisher class, upstream script or unrelated artwork is redistributed.
\par\smallskip\noindent\textit{Reference note.} This article ships no compiled bibliography (biblatex); the entries delivered here are composed from the author's own BibTeX (.bib) fields, with no field invented and no entry dropped.
\begin{defin}\label{deflocalsymmetry}
    A local symmetry acting on a set $\Lambda$ of degrees of freedom is a pair $(G,G_0)$ of finite abelian group $G$ and a generating set $G_0$
    \begin{itemize}
        \item For each $x \in \Lambda$ there is some local symmetry $g \in G$ whose support overlaps $x$.
        \item For each operator $V_X$ which is diagonal in the classical basis and supported in $X$ and each $g \in G$, $g V_X g^{-1}$ is also diagonal in the classical basis and supported in $X$.
        \item The symmetry graph $\mathcal{G}$ is connected.
    \end{itemize}
\end{defin}
The first condition is for simplicity of notation. If there are other degrees of freedom that the local symmetry does not act on, we can simply ignore them in the average. The second is for the simple expression of operator supports. Both of these conditions hold for the symmetry of Wegner gauge theory.

We start with an interaction, expressed as a sum over local interactions $V_X$ supported in subsets $X \subset \Lambda$:
\[\label{eqnlocalinteraction}V = \sum_{X \subset \Lambda} V_X\]
By convention we will take $V_\varnothing = 0$, since this constant does not affect thermodynamic quantities. We want to compute the gauge-averaged effective interaction
\[V_{\rm eff} = -\frac{1}{\beta} \log \mathbb{E}_g (ge^{-\beta V}g^{-1})\]
First we write
\[e^{-\beta V} = \prod_{X \subset \Lambda} e^{-\beta V_X} = \prod_{X \subset \Lambda} (1+f(X,\beta)) = \sum_{n=0}^\infty\Sum_{\{X_1,\ldots,X_n\} \subset 2^\Lambda} \prod_{j=1}^n f(X_j,\beta)\]
where $f(X,\beta) = e^{-\beta V_X}-1$ and $2^\Lambda$ is the power set of $\Lambda$. Linearity of expectation gives
\[e^{-\beta V_{\rm eff}} = \sum_{n=0}^\infty\Sum_{\{X_1,\ldots,X_n\} \subset 2^\Lambda} \mathbb{E}_g\left(g \left(\prod_{j=1}^n f(X_j,\beta)\right) g^{-1}\right).\]

For $X \subset \Lambda$, let $G_0(X)$ be the set of local symmetry generators whose support overlaps $X$. We say two $X,Y \subset \Lambda$ are  ``dependent'' if they have a common overlapping symmetry generator:
\[X \sim Y \Leftrightarrow G_0(X) \cap G_0(Y) \neq \varnothing\]
Since $G$ is finite abelian and generated by $G_0$, the uniform measure on $G$ is the pushforward of the product measure on the abstract product of cyclic groups generated by the elements of $G_0$. Thus functions depending on disjoint subsets of $G_0$ will be independent under gauge averaging. Therefore the expectation value above splits into a product over sets of subsets $\gamma=\{X_1,\ldots,X_m\}$ with connected dependency graphs, which is the graph whose vertices are $\{X_i\}_{i=1,\ldots,m}$ and have edges whenever $X_i \sim X_j$.

We call a set of subsets $\gamma \subset 2^\Lambda$ with a connected dependency graph a ``polymer''. Let $P$ be the set of polymers. We say two polymers $\gamma,\gamma' \in P$ are dependent if
\[\gamma \sim \gamma' \Leftrightarrow \exists X \in \gamma, \exists Y \in \gamma', X \sim Y.\]
(Note that many references, including \cite{kotecky1986cluster}, call this relation ``incompatibility'' and denote it with $\not\sim$.) Otherwise we say they are independent. We can then write
\[e^{-\beta V_{\rm eff}} = \sum_{n=0}^\infty \frac{1}{n!}\Sum_{(\gamma_1,\ldots,\gamma_n) \in P^n \\ \text{mutually independent polymers}} \prod_{j=1}^n \Phi(\gamma_j) \\
\Phi(\gamma) = \mathbb{E}_g \left(g \left(\prod_{X \in \gamma} f(X,\beta)\right) g^{-1}\right)\]
Evaluated in the classical basis, where $H_0$ and $V$ are diagonal, for each fixed configuration of the spins, we can consider this as an abstract polymer model in the language of Koteck\'y and Preiss \cite{kotecky1986cluster}.

Under certain conditions stated in \cite{kotecky1986cluster,fernandez2007cluster}, $V_{\rm eff}$ evaluated on a particular spin configuration has an absolutely convergent cluster expansion. We will apply this cluster expansion pointwise over all spin configurations to get a cluster expansion for $V_{\rm eff}$, while convergence results will be uniform in the spin configurations themselves.

A cluster is an ordered tuple of polymers $C=(\gamma_1,\ldots,\gamma_n) \in P^n$ (may include repeats) with a connected dependency graph, a graph whose vertices are the polymers $\gamma_1,\ldots,\gamma_n$, with an edge whenever $\gamma_i \sim \gamma_j$ are dependent. The cluster expansion is, formally
\[\label{eqnclusterexpapp}-\beta V_{\rm eff} = \sum_{n=1}^\infty \frac{1}{n!} \Sum_{C \in P^n \\ \text{cluster} } \Phi^T(C) \\
\Phi^T(C)=\phi^T(C) \prod_{\gamma \in C} \Phi(\gamma)\]
where $\phi^T(C)=1$ if $n=1$, and otherwise
\[\phi^T(C) = \Sum_{F \subset E(C) \\ F\text{ connected, spanning}} (-1)^{|F|}\]
where $E(C)$ is the set of edges in the dependency graph of $C$, and $F$ is a subset of these edges such that every vertex $\gamma \in C$ has at least one edge in $F$ connected to it and the resulting subgraph is connected. See \cite{fernandez2007cluster} for this explicit formula and a nice review of various convergence criteria.

To place $V_{\rm eff}$ into the general form \eqref{eqnlocalinteraction} of an interaction, we define for a cluster $C = (\gamma_1,\ldots,\gamma_n)$ define the ``support'' of a cluster $C$ as
\[S(C) = \bigcup_{j=1}^n \bigcup_{X \in \gamma_j} X\]
and define
\[\label{eqnlocalinteractioneffective}V_{{\rm eff},X} = - \frac{1}{\beta} \sum_{n \ge 1} \frac{1}{n!} \Sum_{C \in P^n \\ \text{cluster, }S(C)=X} \Phi^T(C).\]
Note that because of the second bullet of Definition \ref{deflocalsymmetry}, this operator has support contained in $X$.

To analyze the convergence of the cluster expansion for gauge averaging, we need several more definitions. As in the main text, we consider the bipartite graph $\mathcal{G}$ with vertex and edge sets
\[V(\mathcal{G}) = \Lambda \cup G_0 \\
E(\mathcal{G})= \{(x,g_0)\ |\ \text{$x$ is in the support of $g_0$}\}\]
For $X \subset \Lambda$ we define
\[G_0(X) := \{\text{set of $g_0 \in G_0$ connected to $X$}\}\]
and for a polymer $\gamma \in P$ we define
\[G_0(\gamma) := G_0(\bigcup_{X \in \gamma} X)\]
We also define two measures of the size of a polymer. The first is the number of gauge generators with support overlapping that of $\gamma$:
\[||\gamma|| := |G_0(\gamma)|= |\bigcup_{X \in \gamma} G_0(X)|\]
To define the second, for each $X \subset \Lambda$ we choose a ``connected hull'', which is a set $H(X) \subset \Lambda$ such that $X \subset H(X)$, $H(X) \cup G_0(H(X))$ is connected in $\mathcal{G}$, and $|G_0(H(X))|$ is minimized over all such sets. Also define
\[\label{eqndiameter}H(\gamma) := \bigcup_{X \in \gamma} H(X) \\
l(X) := |G_0(H(X))| \\
l(\gamma) := |G_0(H(\gamma))|\]
For $\eta \in [0,\infty)$ we define the interaction norm
\[|V|_\eta = \sup_{g_0 \in G_0} \Sum_{X \subset \Lambda \\ g_0 \in G_0(X)} ||V_X||_\infty e^{\eta l(X)},\]
where $||V_X||_\infty$ is the maximum value of $|V_X|$ over all configurations of spins in $X$. This norm is very similar to the one used in \cite{Hastings_2006}, with $l(X)$ serving as an abstract measure of the ``diameter'' of $X$. Intuitively, the $e^{\eta l(X)}$ factor captures the effect that larger weight perturbations are able to generate huge weight operators at a lower order in perturbation theory than smaller weight perturbations, giving the former an exponential boost in destructive power. $\eta^{-1}$ determines a length scale over which the operators in $V$ are required to decay so that $|V|_\eta$ remains finite in the thermodynamic limit.

\subsection{Proof of the Gauge Averaging Theorem}

\begin{thm}\label{apptheoremstatement}
    Let $(G,G_0)$ be a local symmetry satisfying Definition \ref{deflocalsymmetry}. Suppose it has the property that each local generator $g_0 \in G_0$ has support overlapping the supports of at most $\Delta$ other local generators. Then for all $0 < b < \eta-m-\log \Delta - 1$, there is a constant $C_{\rm gauge}(\eta,m,b) > 0$ such that if
    \[\beta |V|_\eta < C_{\rm gauge}(\eta,m,b),\]
    then the cluster expansion is absolutely convergent and gives an exponentially-local interaction \eqref{eqnlocalinteractioneffective} with
    \[\beta |V_{\rm eff}|_m < b.\]
\end{thm}

\begin{proof}
In order to prove convergence, we will consider the abstract polymer model
\[\sum_{n = 0}^\infty \frac{1}{n!} \Sum_{(\gamma_1,\ldots,\gamma_n) \in P^n \\ \text{mutually independent polymers}} \prod_{j=1}^n z(\gamma_j) \\
z(\gamma_j) = ||\Phi(\gamma_j)||_\infty.\]
The absolute cluster expansion of this model dominates the cluster expansion \eqref{eqnclusterexpapp} for every spin configuration. Therefore, absolute convergence for this cluster expansion will prove uniform absolute convergence for \eqref{eqnclusterexpapp}.

Define for each polymer $\gamma \in P$ the ``KP activity''
\[K_{\gamma,b} = \sum_{\gamma' \in P: \gamma' \sim \gamma} ||\Phi(\gamma')||_\infty e^{ml(\gamma') + b||\gamma'||}\]
The ``KP criterion'' is
\[\label{eqnKPcrit}\forall \gamma, K_{\gamma,b} \le b||\gamma|| \qquad \text{KP criterion}\]
Under this condition, absolute convergence of the cluster expansion is guaranteed by the main theorem of \cite{kotecky1986cluster} with $a(\gamma) = b||\gamma||$ and $d(\gamma) = m l(\gamma)$. In the notation of \cite{fernandez2007cluster}, which will be useful below, we have, for $\gamma \in P$
\[\rho_\gamma = ||\Phi(\gamma)||_\infty e^{ml(\gamma)} \\
a_\gamma = b ||\gamma||\]
With this notation \eqref{eqnKPcrit} is Eq. 2.15 of \cite{fernandez2007cluster}.

We want to derive the KP criterion \eqref{eqnKPcrit} from our assumptions. It will be easier to upper bound the ``pinned'' KP activity
\[K_{b,m} = \sup_{g_0 \in G_0} \Sum_{\gamma \in P \\ g_0 \in G_0(\gamma)} ||\Phi(\gamma)||_\infty e^{b||\gamma|| + ml(\gamma)}.\]
First we will show that if $K_{b,m} \le b$, then the KP criterion \eqref{eqnKPcrit} is satisfied. Indeed, if $\gamma' \sim \gamma$, then there is some $g_0 \in G_0(\gamma) \cap G_0(\gamma')$. Thus,
\[K_{\gamma,b} \le \sum_{g_0 \in G_0(\gamma)} \sum_{\gamma': g_0 \in G_0(\gamma')} e^{b||\gamma'||+ml(\gamma')} ||\Phi(\gamma')||_\infty \le ||\gamma|| K_{b,m}.\]
Thus, if we can show
\[\label{eqnpinnedcriterion}K_{b,m} \le b \quad \text{pinned KP criterion}\]
then the KP criterion \eqref{eqnKPcrit} will be satisfied.

Now to upper bound $K_{b,m}$. Let
\[\rho_X = ||f(X,\beta)||_\infty = ||e^{-\beta V_X}-1||_\infty\]
Since $||\Phi(\gamma)||_\infty \le \prod_{X \in \gamma} \rho_X$, we have
\[K_{b,m} \le \sup_{g_0 \in G_0} \Sum_{\gamma \in P \\ g_0 \in G_0(\gamma)} e^{b||\gamma|| + ml(\gamma)} \prod_{X \in \gamma} \rho_X.\]
Note that
\[l(\gamma) \le \sum_{X \in \gamma} l(X) \\
||\gamma|| \le l(\gamma)\]
so for all $\eta \ge 0$,
\[K_{b,m} \le \sup_{g_0 \in G_0} \Sum_{\gamma \in P \\ g_0 \in G_0(\gamma)} e^{-(\eta-b-m) l(\gamma)} \left(\prod_{X \in \gamma} \rho_X e^{\eta l(X)}\right)\]
We now want to group the sum over polymers with a fixed $l(\gamma)$.

Recall this is defined using connected hulls $H(\gamma) \subset \Lambda$ in \eqref{eqndiameter}. For this next bound, we want to use ``connected'' sets of local symmetry generators. Define, for $J \subset G_0$ the set $\Lambda(J)$ of spins in the support of elements in $J$. We will say $J$ is ``connected'' if $J \cup \Lambda(J)$ is connected in $\mathcal{G}$. We define, for each $X \subset \Lambda$ and polymer $\gamma$,
\[J(X) = G_0(H(X)) \\
J(\gamma) = \bigcup_{X \in \gamma} J(X)\]
Note that $J(\gamma)$ is connected since each $J(X)$ is, and since $X \sim Y$ implies $J(X) \cap J(Y) \neq \varnothing$. We have by definition
\[l(\gamma) = |J(\gamma)|.\]
Let
\[\mathcal{J}(l)=\{J \subset G_0\ |\ J\text{ connected and }|J|=l\}.\]
We have
\[K_{b,m} \le \sum_{l=1}^\infty e^{-(\eta-b-m) l} \sup_{g_0 \in G_0}\Sum_{J \in \mathcal{J}(l) \\ g_0 \in J} \Sum_{\gamma \in P \\ J(\gamma)=J} \prod_{X \in \gamma} \rho_X e^{\eta l(X)}\]
Define
\[\mathcal{X}(J) = \{X \neq \varnothing\ |\ J(X) \subset J\}.\]
Polymers $\gamma$ with $J(\gamma) = J$ must be made of sets $X \in \mathcal{X}(J)$. We can thus replace the sum over $\gamma$ with a sum over (nonempty) subsets of $\mathcal{X}(J)$:
\[K_{b,m} \le \sum_{l=1}^\infty e^{-(\eta-b-m) l} \sup_{g_0 \in G_0}\Sum_{J \in \mathcal{J}(l) \\ g_0 \in J} \Sum_{\Gamma \in 2^{\mathcal{X}(J)} \\ \Gamma \neq \varnothing} \prod_{X \in \Gamma} \rho_X e^{\eta l(X)}\]
We can replace the sum of products with an exponential of a sum:
\[\Sum_{\Gamma \in 2^{\mathcal{X}(J)} \\ \Gamma \neq \varnothing} \prod_{X \in \Gamma} \rho_X e^{\eta l(X)} \le \exp\left( \Sum_{X \in \mathcal{X}(J)} \rho_X e^{\eta l(X)} \right) - 1\]
Since $G_0(X) \subset J(X) \subset J$,
\[\sum_{X \in \mathcal{X}(J)}\rho_X e^{\eta l(X)}  \le \sum_{g \in J} \sum_{X: g \in G_0(X)} \rho_X e^{\eta l(X)} \le l q_\eta\]
where we used $l = |J|$ and defined
\[q_\eta = \sup_{g_0 \in G_0} \sum_{X:g_0 \in G_0(X)} \rho_X e^{\eta l(X)}.\]
Note that since $\rho_X \le e^{\beta ||V_X||}-1$,
\[\label{eqnqestimate}q_\eta \le 2 \beta|V|_\eta \quad \text{whenever} \quad \beta |V|_\eta \le \log 2 \]
and in particular goes to zero as $\beta|V|_\eta \to 0$. It will serve as an intermediate measure of the interaction norm.


Putting it all together we have
\[K_{b,m} \le \sum_{l=1}^\infty A(l) e^{-(\eta-b-m)l} (e^{q_\eta l}-1) \\
A(l) = \sup_{g_0 \in G_0} |\{J \in \mathcal{J}(l) : g_0 \in J\}|\]
This bound is true without any assumptions on $\mathcal{G}$. If however $\mathcal{G}$ has bounded symmetry degree $\Delta$, then (see Lemma 9 of \cite{mcdiarmid2021componentstructuredenserandom})
\[A(l) \le (e \Delta)^l\]
Write
\[\alpha = e^{(\log \Delta + 1)-(\eta-b-m)}\]
For whenever
\[\eta-m > \log \Delta + 1\]
then for all $b \in (0,\eta-m-\log \Delta-1)$ we have $\alpha < 1$. In this case, there is some constant $C''(\eta,m,b)$ such that
\[q_\eta < C''(\eta,m,b) \implies \alpha e^{q_\eta} < 1\]
which then implies (summing the above inequality)
\[K_{b,m} \le \frac{\alpha e^{q_\eta}}{1- \alpha e^{q_\eta}} - \frac{\alpha}{1-\alpha}\]
This is a continuous function of $q_\eta$ and goes to zero as $q_\eta \to 0$, therefore there is a another constant $C'(\eta,m,b) \le C''(\eta,m,b)$ such that for all $b \in (0,\eta-m-\log \Delta -1)$,
\[q_\eta < C'(\eta,m,b) \implies K_{b,m} < b.\]
and thus by the argument above \eqref{eqnpinnedcriterion} the KP criterion will be satsified. We will take the constant in the theorem statement to be
\[C_{\rm gauge}(\eta,m,b) = \min(\log 2, C'(\eta,m,b)/2)\]
since by \eqref{eqnqestimate}
\[\beta |V|_\eta < C_{\rm gauge}(\eta,m,b) \implies q_\eta < C'(\eta,m,b)\]


We now wish to show the bound on the interaction norm of $V_{\rm eff}$. This essentially follows the analysis in \cite{fernandez2007cluster}.
We want to compute
\[|V_{\rm eff}|_m=\sup_{g_0 \in G_0} \Sum_{X \subset \Lambda \\ g_0 \in G_0(X)} ||V_{{\rm eff},X}||_\infty e^{m l(X)}\]
The local form of the cluster expansion is given in \eqref{eqnlocalinteractioneffective}. Note that
\[||\sum_{C:S(C)=X} \Phi^T(C)||_\infty \le \sum_{C:S(C)=X} ||\Phi^T(C)||_\infty \\
G_0(S(C)) = \bigcup_{j=1}^n G_0(\gamma_j) \\ l(S(C)) \le \sum_{j=1}^n l(\gamma_j),\]
where the last holds because $\bigcup_{j=1}^n H(\gamma_j) \cup G_0(H(\gamma_j))$ is connected and contains $S(C)$. Putting these together, we have
\[\beta |V_{\rm eff}|_m \le \sup_{g_0 \in G_0} \sum_{n = 1}^\infty \frac{1}{n!} \Sum_{(\gamma_1,\ldots,\gamma_n) \in P^n \\ \text{cluster}} ||\Phi^T(C)||_\infty e^{m \sum_{j=1}^n l(\gamma_j)} \sum_{j=1}^\infty \textbf{1}_{g_0 \in G_0(\gamma_j)} \]
where $\textbf{1}_{g_0 \in G_0(\gamma_j)}$ is a function which is 1 if $g_0 \in G_0(\gamma_j)$ and 0 otherwise. By permutation symmetry of the $\gamma_j$, we can absorb the sum over $j$ of the indicator functions:
\[\beta |V_{\rm eff}|_m \le \sup_{g_0 \in G_0} \sum_{n=1}^\infty \frac{1}{(n-1)!}\Sum_{(\gamma_1,\ldots,\gamma_n) \in P^n \\ \text{cluster, } g_0 \in G_0(\gamma_1)} ||\Phi^T(\gamma_1,\ldots,\gamma_n)||_\infty e^{m\sum_{j=1}^n(\gamma_j)} \]
Furthermore,
\[||\Phi^T(\gamma_1,\ldots,\gamma_n)||_\infty \le |\phi^T(\gamma_1,\ldots,\gamma_n)| \prod_{j=1}^n ||\Phi(\gamma_j)||_\infty\]
So separating off $\gamma_1$, we get
\[\beta |V_{\rm eff}|_m \le \sup_{g_0 \in G_0} \sum_{\gamma_1: g_0 \in G_0(\gamma_1)} ||\Phi(\gamma_1)||_\infty e^{ml(\gamma_1)} \left(1 + \sum_{n=1}^\infty \frac{1}{n!} \sum_{(\gamma_2,\ldots,\gamma_{n+1}) \in P^n} |\phi^T(\gamma_1,\ldots,\gamma_{n+1})| \prod_{j=2}^{n+1} ||\Phi(\gamma_j)||_\infty e^{ml(\gamma_j)} \right)\]
Recalling $\rho_\gamma = ||\Phi(\gamma)||_\infty e^{m l(\gamma)}$, the expression in parentheses is called $\Pi_{\gamma_1}$ in \cite{fernandez2007cluster} Eq. 2.7. They show under the KP criterion that
\[\rho_\gamma\Pi_\gamma \le \rho_\gamma e^{a_\gamma}\]
(see Eq. 2.14 with $\mu_\gamma = \rho_\gamma e^{a_\gamma}$, and also see below Eq. 2.15). This gives
\[\beta |V_{\rm eff}|_m \le \sup_{g_0 \in G_0} \sum_{\gamma: g_0 \in G_0(\gamma)} ||\Phi(\gamma)||_\infty e^{m l(\gamma) + b||\gamma||} = K_{b,m}\]
With our assumptions above we have $K_{b,m} < b$ proving the result.


\end{proof}

\section{Proof of Stability of Wegner Gauge Theory}

Let $\Omega$ be a set of configurations, $P$ the transition matrix of a Markov chain on $\Omega$, and $\mu$ a distribution on $\Omega$. We say $P$ is $\mu$-reversible if it satisfies detailed balance with respect to $\mu$, i.e.
\[\forall \sigma,\sigma' \in \Omega \quad \mu(\sigma) P(\sigma,\sigma') = \mu(\sigma') P(\sigma',\sigma).\]
We define the mixing time as
\[t_{\rm mix} = \inf\{t: \max_{\sigma \in \Omega} ||P^t(\sigma) - \mu||_{TV} \le 1/4\}, \qquad \inf(\varnothing)=\infty\]
where $P^t(\sigma)$ is $P$ applied $t$ times to the initial state $\sigma$ and $||\cdot||_{TV}$ is the total variation distance. The bottleneck lemma says the following:
\begin{lemma}\label{lemmabottleneck} \textbf{Bottleneck lemma:}
    Suppose $A \subset \Omega$ has $0<\mu(A) \le 1/2$ and that there is $\partial A \subset A^c$ such that if $\sigma' \notin A \cup \partial A$, then for all $\sigma \in A$,
    \[\label{eqnbottleneckcondition}P(\sigma,\sigma') = 0.\]
    Then if $P$ is $\mu$-reversible,
    \[t_{\rm mix} \ge \frac14 \frac{\mu(A)}{\mu(\partial A)}\]
\end{lemma}
\begin{proof}
    For $A,B \subset \Omega$ we can consider the conductance
    \[Q(A,B) = \sum_{\sigma \in A, \sigma' \in B} \mu(\sigma)P(\sigma,\sigma').\]
    If every exit from $A$ lands in $\partial A$, then
    \[Q(A,A^c) = Q(A,\partial A).\]
    By $\mu$-reversibility,
    \[Q(A,\partial A) = Q(\partial A,A) \le \mu(\partial A)\]
    By Theorem 7.4 of \cite{levin2026markov}, we thus have
    \[t_{\rm mix} \ge \frac14 \frac{\mu(A)}{\mu(\partial A)}.\]
\end{proof}

\begin{defin}\label{definrangeRmarkov}
    A ball of radius $R$ is a set $B(R,v)$ where $v$ is a fixed lattice vertex, and $B(R,v)$ is the set of edges which lie along a lattice path of length $\le R$ starting at $v$.

    A Markov chain $P$ has range $R$ if for all $\sigma,\sigma' \in \Omega$ with $P(\sigma,\sigma') \neq 0$, there is a $v$ such that $\sigma$ and $\sigma'$ are equal outside $B(R,v)$.
\end{defin}

We need to modify the definition of the bottles $A[C]$ and bottlenecks $\partial A[C]$ from the main text to satisfy the condition \eqref{eqnbottleneckcondition} for a general range $R$ Markov chain. For this, we need a lemma:


\begin{lemma}\label{lemmahomology}\textbf{Homological stability lemma:}
    For every fixed $R$ there are $\varepsilon > 0$, $M$, and $L_0$ such that for all configurations $S$ and $S'$ of Wegner gauge theory on a 3-torus of side length $L > L_0$, which are equal outside of a ball of radius $R$, the following are true:
    \begin{enumerate}
        \item If each component of $\partial S$ has length $< \varepsilon L$, then all but at most $M$ components of $\partial S'$ have length $< \varepsilon L$.
        \item The combined length of these ``long loops'' is less than $L/2$.
        \item The decoded homology classes are equal: $[C(S)] = [C(S')]$.
    \end{enumerate}
\end{lemma}
\begin{proof}
    There is a maximum number $a_0R^3$ of disjoint dual lattice loops (which occupy plaquettes) which neighbor an edge in $B$. We can take $M = a_0 R^3$, then 1 will automatically be satisfied, because only these loops may differ between $\partial S$ and $\partial S'$.
    
    Furthermore, if all loops in $\partial S$ have length $< \varepsilon L$, the total length of such loops neighboring an edge in $B$ must be $\le a_0 R^3 \varepsilon L$. The most their length can be after an update is therefore
    \[\le a_0 R^3 \varepsilon L + b_0 R^3\]
    for another constant $b_0$. For each fixed $R$, we can thus choose $\varepsilon$ independent of $L$ so that for $L > 3b_0 R^3$, this is less than $L/2$.

    For stability of the decoded homology class, let $\gamma_i$, $i=1,\ldots, k \le M$ denote the set of components of $\partial S$ neighboring an edge in $B$, and let $\gamma_i'$, $i=1,\ldots,k' \le M$ denote the set of components of $\partial S'$ neighboring an edge in $B$. All other components of $\partial S$ and $\partial S'$ are the same, so the difference of the decoded 2-cycles is
    \[C(S)+C(S')= S + S'+ \sum_{i=1}^k D(\gamma_i) + \sum_{j=1}^{k'} D(\gamma_j').\]
    $S+S'$ is supported in $B$ because that's where the update occurs. Furthermore, each $D(\gamma_i)$ or $D(\gamma_j')$ is supported in a bounding cube of side length $<L/4$, and all of these cubes intersect the ball of radius $R$. Their union is thus contained in a cube of side length $L/2 + 2R+2$, where the $+2$ accounts for the shift between the lattice and its dual. This is contractible as long as $L > 4R+4$. So let us choose
    \[L_0 = \max(4R+4,3b_0R^3)\]
    In this case, $C(S) + C(S')$ must be a boundary, so point 3 is satisfied.
\end{proof}

With this lemma in mind, we define, for each $R$, $A[C]$ to be the set of configurations $S$ with $\partial S$ having all components of length $< \varepsilon L$ (call these ``short''), and $S$ having decoded homology class $[C]$, and $\partial A[C]$ to be the set of $S$ where all but $k \in [1,M]$ components of $\partial S$ are short, and the rest have length at least $\varepsilon L$, but with total length less than $L/2$, and having decoded homology class $[C]$. By Lemma \ref{lemmahomology}, any range $R$ Markov chain must pass through $\partial A[C]$ to leave $A[C]$.



\begin{thm}\label{thmWegnermemoryapp} \textbf{Stability of Self-Correcting Memory:}
    Let $\Lambda^L$ denote the $L \times L \times L$ cubic lattice with periodic boundary conditions, and $H_0^L$ the Wegner gauge theory on this lattice. Let $\beta = 1/T$. For all
    \[0 < T < 1/\log 5 \qquad \eta > \log 6 + 1\]
    there exists $C_{\rm stab}(T,\eta) > 0$ with the following property. For every $R$ there exist constants
    \[A_0(T,\eta,R) > 0, \qquad \lambda(T,\eta,R) > 0, \qquad L_0(T,\eta,R)\]
    such that for every $L \ge L_0$, every interaction $V$ on $\Lambda^L$ satisfying
    \[\beta |V|_\eta < C_{\rm stab}(T,\eta)\]
    and every range $R$ Markov chain on $\Lambda^L$ which is $\mu$-reversible, where $\mu \propto \exp(-\beta(H_0^L+V))$, its mixing time satisfies
    \[t_{\rm mix} > A_0 e^{\lambda L}\]
\end{thm}

\begin{proof}
For fixed $R$, $L_0$ and $\varepsilon$ are defined in Lemma \ref{lemmahomology}, and the bottles $A[C]$ and bottlenecks $\partial A[C]$ below that. Our goal is to show for the perturbed Gibbs measure $\frac{1}{Z} e^{-\beta(H_0 + V)}$ that the bottleneck ratio
\[\mu(\partial A[C])/\mu(A[C])\]
is upperbounded by $e^{-\Omega(L)}$ as $L \to \infty$. By Lemma \ref{lemmabottleneck}, this will show an exponentially long mixing time for any Markov dynamics satisfying detailed balance for $\mu$.

Since $A[C]$ and $\partial A[C]$ are $G$-symmetric sets, we have
\[\mu(\partial A[C])/\mu(A[C]) = \mu^G(\partial A[C])/\mu^G(A[C]),\]
where $\mu^G$ is the gauge averaged Gibbs measure $\frac{1}{Z} e^{-\beta(H_0 + V_{\rm eff})}$. Thus it will suffice to upperbound the bottleneck ratio on the right-hand side.

As in the main text, we have an injective map
\[i:\partial A[C] \hookrightarrow A[C] \times \Theta\]
where $\Theta$ is the set of size-$k \in [1,M]$ sets of loops of length in $[\epsilon L,L/2)$ whose total length is less than $L/2$. In particular, for $\Gamma \in \Theta$ define
\[D(\Gamma) = \sum_{\gamma \in \Gamma} D(\gamma).\]
Then $i(S) = (S + D(\Gamma),\Gamma)$. Also set
\[|\Gamma| = \sum_{\gamma \in \Gamma} |\gamma|.\]
We have
\[\label{eqnpeierlsstepone}\mu^G(\partial A[C]) = \frac{1}{Z} \sum_{(S',\Gamma) \in \text{Im}(i)} e^{-\beta(|\partial S'|+|\Gamma|+V_{\rm eff}(S'+D(\Gamma)))} \\
\le \frac{1}{Z} \sum_{S' \in A[C]} e^{-\beta(|\partial S'| + V_{\rm eff}(S'))} \sum_{\Gamma \in \Theta} e^{-\beta(|\Gamma| +  V_{\rm eff}(S'+D(\Gamma))-V_{\rm eff}(S'))}\]
In order to run the Peierls bound, we desire an upper bound on the size of
\[\delta_\Gamma V_{\rm eff}(S) :=V_{\rm eff}(S+D(\Gamma)) - V_{\rm eff}(S)\]
For each $X \subset \Lambda$, $V_{{\rm eff},X}$ in \eqref{eqnlocalinteractioneffective} is gauge invariant. So for it to contribute to $\delta_\Gamma V_{\rm eff}(S)$, it must contain such a lattice 1-cycle $C'$ linking $\Gamma$. In this case, its contribution is upper bounded by $2||V_{{\rm eff},X}||_\infty$. Therefore,
\[|\delta_\Gamma V_{\rm eff}(S)| \le 2\Sum_{X \subset \Lambda \\ \exists C' \subset X , C'\text{ linking }\Gamma} ||V_{{\rm eff},X}||_\infty \]
The linking cycle $C'$ must either be contractible or have length at least $L$ to wrap the 3-torus. Let us therefore separate the sum into ``small'' sets $X$ with $l(X) < L$ and ``large'' sets $X$ with $l(X) \ge L$. The contribution from large sets may be upper bounded by
\[2\sum_{g_0 \in G_0} \Sum_{X:g_0 \in G_0(X) \\ l(X) \ge L} ||V_{{\rm eff},X}||_\infty \le 2L^3 \sup_{g_0 \in G_0} \Sum_{X:g_0 \in G_0(X) \\ l(X) \ge L} ||V_{{\rm eff},X}||_\infty e^{m(l(X)-L)} \\
\le 2L^3 e^{-mL} |V_{\rm eff}|_m =: R_m\]
As long as $|V_{\rm eff}|_m < \infty$, which we will guarantee below, $R_m$ will be an exponentially small remainder which does not contribute to the large $L$ Peierls estimate.


Now we proceed to bound the contribution from small sets $X$ containing a contractible loop $C'$ linking $\Gamma$. Suppose $C'$ has length $r$. Then, since $C'$ is contractible, it is contained in a bounding cube of side length $r/2$. Since it links $\Gamma$, $\Gamma$ must pass through this cube as well. Thus, there is a lattice path $p$ of length $l(p) \le (3/2)r$ from a vertex of $C'$ to a vertex at a corner of a plaquette of $\Gamma$. Let $N(\Gamma,(3/2)r) \subset G_0$ denote the set of vertices which can be connected to $\Gamma$ by such a path. By the above, we have $G_0(C') \cap N(\Gamma,(3/2)r) \neq \varnothing$. Thus, we can give a bound
\[\Sum_{X \subset \Lambda \\ \exists C' \subset X \\ C'\text{ linking }\Gamma, \text{ contractible}} ||V_{{\rm eff},X}||_\infty \le \sum_{r=4}^{L-1} \sum_{g_0 \in N(\Gamma,(3/2)r)} \Sum_{X: g_0 \in G_0(X) \\
l(X) \ge r} ||V_{{\rm eff},X}||_\infty \\
\le |V_{\rm eff}|_m\sum_{r=4}^{L-1} |N(\Gamma,(3/2)r)| e^{-mr}\]
$|N(\Gamma,(3/2)r)|$ is the number of points in an $\ell^1$-metric tube around $\Gamma$. Thus there is a universal constant $c_0$ such that
\[|N(\Gamma,(3/2)r)| \le c_0 r^3 |\Gamma|.\]
Let
\[C_{\rm perim}(m) = 2c_0 \sum_{r=4}^{\infty} r^3 e^{-mr}\]
This is finite and goes to 0 as $m \to \infty$. We have
\[|\delta_\Gamma V_{\rm eff}(S)| \le C_{\rm perim}(m) |V_{\rm eff}|_m|\Gamma| + R_m\]

Combined with \eqref{eqnpeierlsstepone} we obtain
\[\mu^G(\partial A[C]) \le \frac{e^{\beta R_m}}{Z} \sum_{S' \in A[C]} e^{- \beta (|\partial S'|+V_{\rm eff}(S'))}\sum_{\Gamma \in \Theta} e^{-\beta(1 -  C_{\rm perim}(m)|V_{\rm eff}|_m) |\Gamma|} \\
= e^{\beta R_m} \mu(A[C])\sum_{\Gamma \in \Theta} e^{-\beta(1 -  C_{\rm perim}(m)|V_{\rm eff}|_m) |\Gamma|}.\]
We see at this stage that $V_{\rm eff}$ has just served to renormalize the tension of the long loops $\Gamma$. We need to ensure that the tension is still within the radius of convergence of the Peierls estimate. Let
\[\label{eqnqparam} l_* = \lceil \varepsilon L\rceil \qquad q = 5 e^{-\beta(1-C_{\rm perim}(m)|V_{\rm eff}|_m)}.\]
Let $\Theta_k$ be the subset of $\Theta$ consisting of those $\Gamma$ with $k$ components. We have, as long as $q < 1$,
\[\sum_{\Gamma \in \Theta} e^{-\beta( +  C_{\rm perim}(m)|V_{\rm eff}|_m) |\Gamma|} = \sum_{k=1}^M \sum_{\Gamma \in \Theta_k} e^{-\beta( +  C_{\rm perim}(m)|V_{\rm eff}|_m) |\Gamma|} \\
\le \sum_{k=1}^M \left(L^3\sum_{r=l_*}^\infty q^r\right)^k
\\ = \sum_{k=1}^M \left(\frac{L^3 q^{l_*}}{1-q} \right)^k\]
For all $q < 1$, this function has exponential decay $O(\exp\left(-\lambda L \right))$ for all
\[0<\lambda < \varepsilon (-\log q).\]
We will choose $C_{\rm stab}(T,\eta)$ so that for all $V$, $q$ will be bounded away from 1 uniformly in $L$, it will suffice for the theorem to choose $\lambda = \varepsilon (-\log q)/2$ and $A_0$ an appropriately large constant. Indeed, this gives an upper bound on the bottleneck ratio $\mu^G(\partial A[C])/\mu^G(A[C])$ which goes to zero exponentially quickly with $L$. Since on $T^3$, there are 8 possible $[C]$, and all the $A[C]$ are disjoint, one of them will have $\mu^G(A[C]) < 1/2$. Thus, by Lemma \ref{lemmabottleneck}, we will get the desired bound on the mixing time.

To finish the proof, suppose that $T < 1/\log 5$, i.e.
\[\beta > \log 5.\]
Then for all $\eta > \log 6 + 1$, we can choose $m$ and $b$ such that
\[
0 < m < \eta - \log 6 -1 \\
0 < b < \min\left(\eta-m - \log 6 - 1, \frac{\beta -\log 5}{C_{\rm perim}(m)}\right)\]
By the Gauge Averaging Theorem (Theorem \ref{apptheoremstatement}) (here $\Delta = 6$), so long as
\[\beta |V|_\eta < C_{\rm gauge}(\eta,m,b),\]
we obtain a gauge-invariant effective interaction $V_{\rm eff}$ with $\beta|V_{\rm eff}|_m < b$. In particular, by the condition on $b$, $q$ in \eqref{eqnqparam} will be bounded uniformly in $L$ by
\[q < 5e^{-\beta+b C_{\rm perim}(m)} < 1\]
Thus, the theorem is proved with $C_{\rm stab}(T,\eta) = C_{\rm gauge}(\eta,m,b)$.

\end{proof}


\begin{thebibliography}{99}
\bibitem[]{Hastings_2006} Hastings, Matthew B.; Koma, Tohru. Spectral Gap and Exponential Decay of Correlations. Communications in Mathematical Physics 265 (2006) 781–804
\bibitem[]{fernandez2007cluster} Fern{\'a}ndez, Roberto; Procacci, Aldo. Cluster expansion for abstract polymer models. New bounds from an old approach. Communications in Mathematical Physics 274 (2007) 123--140
\bibitem[]{kotecky1986cluster} Koteck{\`y}, Roman; Preiss, David. Cluster expansion for abstract polymer models. Communications in Mathematical Physics 103 (1986) 491--498
\bibitem[]{levin2026markov} Levin, David A; Peres, Yuval. Markov chains and mixing times. (2026)
\bibitem[]{mcdiarmid2021componentstructuredenserandom} Colin McDiarmid; Alex Scott; Paul Withers. The component structure of dense random subgraphs of the hypercube. (2021)
\end{thebibliography}
\end{document}
